\section{局域交换：Slater 势与优化有效势}

精确 HF 交换算符非局域，计算昂贵。GV §2.8 给出两类局域化方案，是 LDA 交换的前身。

\subsection{Slater 局域交换}

对每个占据轨道定义轨道依赖交换势 $V_{x,\alpha}(\mathbf{r})$，再按密度加权平均：
\begin{equation}
  V_x^{S}(\mathbf{r})
  =\frac{
  \sum_\alpha n_\alpha|\phi_\alpha(\mathbf{r})|^2 V_{x,\alpha}(\mathbf{r})
  }{
  \sum_\alpha n_\alpha|\phi_\alpha(\mathbf{r})|^2
  }.
\label{eq:slaterVx}
\end{equation}
利用 Slater 行列式的 $g$ \eqref{eq:g_slater}，可证
\begin{equation}
  V_x^{S}(\mathbf{r})
  =\int\mathrm{d}\mathbf{r}'\,
  \frac{e^2 n(\mathbf{r}')}{|\mathbf{r}-\mathbf{r}'|}
  \bigl[g(\mathbf{r},\mathbf{r}')-1\bigr].
\end{equation}
即 Slater 势是\textbf{交换孔的静电势}。均匀气中它是常数，等于 $2\varepsilon_x$（每自旋平均交换能的两倍），渐近 $r\to\infty$ 时有限体系 $V_x^S\sim -e^2/r$。

\subsection{优化有效势（OEP）}

要求局域势 $V_x(\mathbf{r})$ 使由该势产生的 Slater 行列式能量最低，得到积分方程（Sharp--Horton / Talman--Shadwick）。OEP 比 Slater 更接近精确 HF，且自然给出正确的 $-e^2/r$ 渐近。实践中常用 KLI 近似求解。

\subsection{与 LDA 交换的关系}

均匀气 $\varepsilon_x(n)=-B_d n^{1/d}\times\text{const}$，其函数导数
\begin{equation}
  v_x^{\mathrm{LDA}}(n)
  =\frac{\mathrm{d}(n\varepsilon_x)}{\mathrm{d}n}
  =\frac{d+1}{d}\varepsilon_x(n)
\end{equation}
即 Dirac/Gaspar 势。三维
\begin{equation}
  v_x^{\mathrm{LDA}}
  =-\frac{e^2}{\pi}(3\pi^2 n)^{1/3}
  =-\frac{e^2 k_F}{\pi},
\end{equation}
恰为 $\Sigma^{\mathrm{x}}(k_F)$，而非 Slater 的 $2\varepsilon_x=-(3/2)e^2k_F/\pi$。KS 方程应使用函数导数而非 Slater 平均——这是 DFT 与 HF 局域化之间的关键差别。
